This appendix describes the computational support for the manuscript. The code
and notebooks in \texttt{code/} are provided to visualize and explore the
proposed spectral-coherence structure numerically. Computational results are
illustrative and do not substitute for formal proofs.

\subsection{What the Code Does}

\begin{itemize}
  \item \texttt{code/notebooks/spectral-visualization.ipynb}: Plots the magnitude
        and phase of $\zeta(s)$ along and near the critical line. Visualizes the
        reflection symmetry $s \mapsto 1-s$.
  \item \texttt{code/notebooks/critical-line-geometry.ipynb}: Plots the critical
        line and the placement of known nontrivial zeros. Illustrates the geometric
        intuition of the balance framework.
  \item \texttt{code/src/placeholder.py}: Placeholder module with utility functions
        for computing $\zeta(s)$ numerically (using the \texttt{mpmath} library
        for high-precision arithmetic).
\end{itemize}

\subsection{Numerical Limitations}

All numerical computations of $\zeta(s)$ are approximate. The known zeros have
been computed to many decimal places by other researchers (see \cite{Odlyzko1987,
PlattTrudgian2021}), and the code in this repository does not aim to extend the
numerical verification of RH. Rather, it supports visual understanding of the
framework.

No claim in the manuscript relies on the output of these computations.

\subsection{Reproducing the Computations}

To reproduce the notebooks:

\begin{enumerate}
  \item Install Python 3.9 or later with the packages listed in
        \texttt{code/README.md}.
  \item Navigate to \texttt{code/notebooks/} and launch Jupyter:
        \texttt{jupyter notebook}.
  \item Run all cells in order.
\end{enumerate}

See \texttt{code/README.md} for full instructions.
